% 核心结论：两条计算路线都给出2z。
% 图契约：物理竞赛讲义；印刷宽16cm，最小字号8.5pt；几何示意与公式分区。
% 布局：[几何/曲线 7cm] [定义/结果 7cm]；实线主量，虚线辅助；无装饰数据。
\documentclass[tikz,border=6pt]{standalone}
\usepackage{fontspec,xeCJK,amsmath,amssymb,graphicx}
\IfFontExistsTF{Times New Roman}{\setmainfont{Times New Roman}}{\setmainfont{TeX Gyre Termes}}
\IfFontExistsTF{Heiti SC}{\setCJKmainfont[AutoFakeBold=1.5]{Heiti SC}}{\setCJKmainfont[AutoFakeBold=1.5]{FandolHei-Regular}}
\usetikzlibrary{arrows.meta,bending,calc,positioning,patterns}
\definecolor{ink}{HTML}{243745}
\definecolor{blue}{HTML}{275DAB}
\definecolor{green}{HTML}{177C63}
\definecolor{red}{HTML}{BD4935}
\definecolor{orange}{HTML}{AA6C16}
\definecolor{purple}{HTML}{715393}
\tikzset{
 every node/.style={font=\fontsize{9.5}{13}\selectfont,text=ink},
 axis/.style={-{Stealth[length=4pt]},line width=.65pt,draw=ink!65},
 vec/.style={-{Stealth[length=5pt]},line width=1pt,draw=blue},
 short/.style={-{Stealth[length=3pt]},line width=.8pt},
 guide/.style={dashed,draw=ink!40,line width=.5pt},
 card/.style={rounded corners=2pt,draw=blue!60,fill=blue!4,
   align=center,text width=6.8cm,inner xsep=9pt,inner ysep=7pt},
 title/.style={font=\fontsize{11}{14}\selectfont\bfseries,text width=15cm,align=center},
 label/.style={font=\fontsize{8.5}{11}\selectfont,align=center}}
\begin{document}
\begin{tikzpicture}[x=1cm,y=1cm]
\path[use as bounding box] (-8,-3.35) rectangle (8,3.35);
\node[title] at (0,2.95) {例：用直接计算核对叉乘场的散度};
\begin{scope}[shift={(-4,0)}]

\draw[axis] (-2.4,0)--(2.5,0) node[right] {$x$};\draw[axis] (0,-2)--(0,2.2) node[above] {$y$};
\foreach \a in {0,60,...,300}{\draw[short,blue] (\a:1.3) arc[start angle=\a,end angle=\a+30,radius=1.3];}
\node[label] at (0,-2.45) {$\mathbf A=(-y,x,0)$，$\nabla\times\mathbf A=2\mathbf e_z$};

\end{scope}\node[card,draw=blue!65,fill=blue!4] at (3.9,1.65) {$\mathbf B=(0,0,z)$，$\nabla\times\mathbf B=\mathbf0$\\$\mathbf A\times\mathbf B=(xz,yz,0)$};
\node[card,draw=green!65,fill=green!4] at (3.9,0) {直接求导：\\$\nabla\cdot(\mathbf A\times\mathbf B)=z+z=2z$};
\node[card,draw=orange!65,fill=orange!4] at (3.9,-1.65) {恒等式：\\$\mathbf B\cdot(\nabla\times\mathbf A)-\mathbf A\cdot(\nabla\times\mathbf B)=2z$};
\end{tikzpicture}\end{document}
